\documentclass[12pt]{article}
\usepackage[utf8]{inputenc}
\usepackage[T1]{fontenc}
\usepackage{hyperref}
\title{Gowlan Language Resources for K-12 Curriculum Development}
\author{Compiled by Claude Sonnet 5.5 (claude-sonnet-5-5) for Andr\'as Kornai}
\date{}
\begin{document}
\maketitle

\section{Introduction}
Gowlan (ISO 639-3 goj) is listed in the reference catalogues as a small language of central India (Maharashtra, Amravati District), with a population of about 15,000 according to Global Recordings Network. Its status is doubtful: Glottolog treats the code as a bookkeeping entry and its reference for the Gowari caste treats Gowlan as a caste name rather than a distinct language. It should not be confused with Gojri (gju). My searches found no curriculum material of any kind for this language.

\section{Primary Educational Resources}
None found. No primers, textbooks or school programmes for Gowlan appeared in search results.

\section{Linguistic and Reference Materials}
No grammar or dictionary was found. The only reference-level information located is the catalogue record of Global Recordings Network: \url{https://globalrecordings.net/en/language/goj}.

\section{Religious and Community Resources}
None found for Gowlan specifically.

\section{Digital and Audio Resources}
The Global Recordings Network language page (\url{https://globalrecordings.net/en/language/goj}) is the only audio-related listing that surfaced; its content was not examined beyond the search result.

\section{Curriculum Development Context}
Because the existence of Gowlan as a language distinct from surrounding languages (for example Korku and regional Indo-Aryan varieties) is itself in question, no native-language curriculum can responsibly be planned without first field-verifying whether a speech community with a separate variety exists. Any school instruction in this area would in practice be in Marathi or Hindi.

\end{document}